% =====================================================================
% SECTION I — INTRODUCTION
% File: section1_introduction.tex
% =====================================================================

% ------------------------------------------------------------------
\section{Introduction}
\label{sec:intro}
% ------------------------------------------------------------------

Reliable demand forecasting for mineral commodities and industrial
supply chains is a critical enabler of procurement planning,
resource allocation, and geopolitical risk management. Commodity
demand is jointly shaped by two fundamentally different signals:
\textit{temporal} patterns arising from historical consumption
trends and seasonality, and \textit{structural} patterns arising
from inter-commodity dependencies encoded in supply-chain graphs.
State-of-the-art forecasting models address these signals in
isolation: transformer-based sequence models~\cite{lim2021tft}
exploit temporal dynamics but ignore graph connectivity; graph neural
networks~\cite{velickovic2018gat,wu2021gnn} exploit structural
relationships but require fixed temporal aggregation. Hybrid models
that combine both branches~\cite{ding2023mstgat} apply a
\textit{fixed} pre-specified weight to each modality, assuming a
universal balance between temporal and structural signal that does
not generalise across datasets or forecasting regimes.

We observe that this assumption is empirically wrong. Daily
fast-moving consumer goods (FMCG) data are dominated by temporal
dynamics with rich intra-distribution history; annual mineral
commodity data exhibit ambiguous modality dominance because
structural supply-chain dependencies are as informative as
sparse historical trends. A model that fixes the temporal-to-structural
ratio cannot simultaneously serve both regimes optimally.

This paper proposes the \textbf{Adaptive Temporal-Structural Fusion
(ATSF)} model, which learns a data-driven fusion weight $\alpha$,
conditioned on material type and forecast horizon, that determines
the optimal balance between temporal and structural representations
on a per-instance basis. The resulting scalar weight is fully
interpretable: post-training, it directly quantifies each dataset's
reliance on temporal vs.\ structural signal.

Our principal contributions are:

\begin{enumerate}[leftmargin=*, nosep]
  \item A novel adaptive fusion mechanism conditioned on material
    type and forecast horizon that replaces fixed-weight blending
    and entropy regularisation with one-sided boundary
    regularisation, enabling strong unimodal preference when
    empirically supported.
  \item Empirical demonstration that the fusion weight correctly
    identifies the dominant modality across two structurally
    distinct datasets: $\alpha{=}0.919{\pm}0.010$ on daily
    SupplyGraph (temporal dominance) and $\alpha{=}0.489{\pm}0.197$
    on annual USGS mineral data (mixed regime), a 43-percentage-point
    structural contribution gap that constitutes the paper's core
    finding.
  \item State-of-the-art performance: 30.1\% WAPE reduction over the
    best neural alternative on SupplyGraph, and a 95\% confidence
    interval for USGS WAPE entirely below the strongest classical
    baseline.
  \item A probabilistic uncertainty quantification layer with
    calibrated P10--P90 prediction intervals and a rule-based
    risk-scoring interface for procurement decision support.
\end{enumerate}

\smallskip
\noindent\textbf{Disclosure.} This work uses two publicly available
datasets: SupplyGraph~\cite{wasi2024supplygraph} and USGS Mineral
Resources~\cite{usgs2023minerals}. No human subjects, protected
health information, or personally identifiable data are involved.
Experiments were conducted on Google Cloud Platform~(Vertex AI
Workbench); systematic recording of GPU-hour usage and associated
CO\textsubscript{2} emissions was not performed.
